\documentclass[10pt,a4paper]{amsart}
\usepackage[margin=19mm]{geometry}
\usepackage{amsmath,amssymb,mathtools}
\usepackage[scr=rsfs,cal=euler]{mathalfa}
\usepackage{xfrac}
\usepackage[unicode,hidelinks,bookmarks=false]{hyperref}
\newcommand{\ie}{i.e.\ }
\newcommand{\cf}{cf.\ }
\newcommand{\resp}{respectively\ }
\newcommand{\Subsection}{\subsection}
\providecommand{\nobreakdash}{\nobreak\hbox{-}}
\setlength{\emergencystretch}{2em}
\multlinegap0pt
\usepackage{amssymb}
\usepackage{mathtools}
\usepackage{xcolor}
\usepackage{dsfont}
\usepackage{enumerate}
\usepackage{float}
\newtheorem{lemma}{Lemma}
\def\H{\mathbb H}
\def\P{\mathbb P}
\title{Classification of Quaternionic Projective Transformations by Equicontinuity Regions}
\author{Sandipan Dutta, Krishnendu Gongopadhyay, Rahul Mondal}
\date{Source: arXiv:2511.20053v2 (CC BY 4.0)}
\begin{document}
\maketitle

\noindent\emph{Source and licence.} Classification of Quaternionic Projective Transformations by Equicontinuity Regions. Version arXiv:2511.20053v2 (CC BY 4.0), distributed by arXiv under the Creative Commons Attribution 4.0 International licence, http://creativecommons.org/licenses/by/4.0/, as recorded in the arXiv licence metadata for this exact version and shown on the abstract page https://arxiv.org/abs/2511.20053v2. Full licence notice: \texttt{licenses/arxiv-2511.20053-CC-BY-4.0.txt}.\par\smallskip

\noindent\textit{Editorial scope.} Counted unit: the article's lemma th:par4 (source0.tex lines 719--734). The article's complete original proof of it is delivered with it (source0.tex lines 736--791, one \texttt{proof} environment of 56 lines). No mathematics is written, repaired or re-derived: the statement and the proof are the article's own lines, in the article's own notation, and no retained line is substituted or re-flowed.  No further source-local context is needed: every cross-reference of every retained line resolves inside the two delivered spans.  Nothing was omitted from any retained span, so the package carries no omission record.  No retained line contains a citation, so the wrapper carries no bibliography and no bibliography entry.  No retained line contains a figure environment, an image-inclusion command or drawing code, so no image is delivered and the package declares no asset dependency.  Editorial formatting only, in addition: an AMS \texttt{amsart} preamble that restates the article's own declarations and macros used by the retained lines, namely the environments \texttt{lemma} on separate counters, together with the standard packages those lines need.  The retained environments print in this document as Lemma 1 in order of appearance, whereas the article prints the same items with the numbering of its own full document.  The complete native source of the article as distributed in the arXiv e-print for this exact version is bundled unchanged as \texttt{sources/189758/source0.tex}.
\begin{lemma}
\label{th:par4}
Let $\tilde{\gamma} \in \mathrm{PSL}(n+1, \mathbb{H})$ be a parabolic transformation whose lift is given by 
\[
\gamma =
\begin{bmatrix}
    \mathrm{J}(\lambda, k) & \\
    & \mathrm{J}(\lambda, l)
\end{bmatrix}
\in \mathrm{SL}(n+1, \mathbb{H}), \qquad k + l = n + 1, \; k < l.
\]
Then the equicontinuity region of the cyclic group $\Gamma := \langle \tilde{\gamma} \rangle$ is 
\[
\mathrm{Eq}(\Gamma) = \P^n_\H\setminus L\{ e_1, e_2, \ldots, e_{k+l-1} \}.
\]
\end{lemma}

\begin{proof}
The $m$-th power of $\gamma$ is given by
\[
\gamma^m =
\begin{bmatrix}
    \mathrm{J}(\lambda, k)^m & \\
    & \mathrm{J}(\lambda, l)^m
\end{bmatrix}.
\]
Consider the sequence
\[
\frac{1}{{m \choose l - 1}} \gamma^m.
\]
Since \(k < l\), the ratio of binomial coefficients satisfies;
\[
\frac{\binom{m}{\,k-1\,}}{\binom{m}{\,l-1\,}} \longrightarrow 0 
\qquad \text{as } m \to \infty.
\]

and 
\[
\frac{1}{{m \choose l - 1}} \mathrm{J}(\lambda, l)^m \longrightarrow a\,\mathrm{E}_{1,l}
\]
for some $|a| = 1$, it follows that
\[
\frac{1}{{m \choose l - 1}} \gamma^m \longrightarrow
g =
\begin{bmatrix}
    \mathbf{0} & \\
    & a\,\mathrm{E}_{1,l}
\end{bmatrix}  \in \mathrm{QP}(n+1,\mathbb{H})
\]
Hence,
$\ker(g) = L\{ e_1, e_2, \ldots, e_{k+l-1} \}.$


Similarly, for the inverse powers $\gamma^{-m}$, the sequence
$\frac{1}{{m + l - 2 \choose l - 1}} \gamma^{-m}$
converges to the same limit
\[
g =
\begin{bmatrix}
    \mathbf{0} & \\
    & a\,\mathrm{E}_{1,l}
\end{bmatrix} \in \mathrm{QP}(n+1,\mathbb{H})
\]
and therefore,
$\ker(g) = L\{ e_1, e_2, \ldots, e_{k+l-1} \}.$


Combining the two cases, we conclude that
\[
\mathrm{Eq}(\Gamma) =  \P^n_\H\setminus L\{ e_1, e_2, \ldots, e_{k+l-1} \} .
\]
This completes the proof. 
\end{proof}

\end{document}
